\documentclass[12pt]{article}

\usepackage[portuguese]{babel}
\usepackage{float}
\usepackage{listings}
\usepackage{indentfirst}
\usepackage{amsfonts, amssymb, amsmath}
\usepackage{graphicx}
\usepackage{enumitem}
\usepackage[colorlinks=true, allcolors=blue]{hyperref}
\usepackage[letterpaper,
            top=1.5cm,
            bottom=1.5cm,
            left=2.75cm,
            right=2.75cm,
            marginparwidth=1.75cm]{geometry}

\title{Lista Revisão -- Prova 1}
\author{Guilherme Lopes}
\date{\today}

\begin{document}
\maketitle

\begin{enumerate}
    \item Defina formalmente o conceito de alfabeto ($\Sigma$) e símbolo. Em seguida, 
    forneça um exemplo de um alfabeto onde pelo menos um dos símbolos seja composto 
    por mais de um caractere.
    \item Considere o alfabeto $\Sigma = \{a, b, c, de\}$ . De acordo com a definição 
    de comprimento $|x|$, calcule o tamanho das seguintes cadeias, justificando se 
    pertencem ou não a $\Sigma^$ :
    \begin{enumerate}[label=\alph*)]
        \item $x = ababac$
        \item $y = abdec$
        \item $z = d$
    \end{enumerate}
    \item Defina o conceito de Linguagem sobre um alfabeto $\Sigma$ . Forneça um 
    exemplo de uma linguagem finita e uma linguagem infinita (utilizando a notação 
    de conjuntos).
    \item Explique a diferença técnica entre o Fechamento de Kleene ( $\Sigma^$ ) 
    e o Fechamento Positivo ( $\Sigma^+$ ). Sob qual condição fundamental o conjunto 
    $\Sigma^*$  é considerado infinito? Dica Didática: Lembre-se da relação de união 
    de potências do alfabeto.
    \item Dada a cadeia $w = abcde$ , liste exaustivamente todos os prefixos e todos 
    os sufixos possíveis desta cadeia. Não esqueça do elemento neutro $\lambda$ .
    \item Dados $L_1 = \{a, ab\}$ e $L_2 = \{\lambda, a, ba\}$, calcule:
    \begin{enumerate}[label=\alph*)]
        \item $L_1 \cup L_2$
        \item $L_1 \cap L_2$
    \end{enumerate}
    \item Um Autômato Finito Determinístico é definido pela quíntupla $\langle 
        \Sigma, S, S_0, \delta, F \rangle$. Explique o que cada elemento representa. 
        Em seguida, forneça a quíntupla específica que descreveria uma linguagem 
        que contenha apenas a união resultante do Exercício 06 ( $L_1 \cup L_2$ ), 
        assumindo um estado final para cada cadeia.
\end{enumerate}

2. Nível Médio: Operações, AFD e AFND
Exercício 08 (Concatenação de Linguagens):Utilizando as linguagens  $L_1 = \{0, 11\}$  e  $L_2 = \{0, 1, 00\}$ , determine o conjunto resultante da operação de concatenação  $L_1.L_2$ .Exercício 09 (Modelagem de Problema Real):Com base no exemplo do micro-ondas (Slide 02), modele um sistema de  porta eletrônica inteligente . Identifique:a) O conjunto de entradas (ex: senha_correta, sensor_abertura).b) O conjunto de estados (incluindo: Trancada, Processando_Senha, Aberta).Exercício 10 (Análise de Tabela de Transição):Considere o AFD definido pela tabela abaixo, com  $\Sigma = \{0, 1\}$ ,  $S = \{S_0, S_1\}$ , estado inicial  $S_0$  e  $F = \{S_1\}$ .| Estado Atual | Símbolo 0 | Símbolo 1 || ------ | ------ | ------ || $S_0$ | $S_1$ | $S_0$ || $S_1$ | $S_1$ | $S_0$ |
Descreva, em linguagem natural, qual o padrão de cadeias aceito por este dispositivo.Exercício 11 (Problema HLCR):No clássico problema do Homem, Lobo, Cabra e Repolho, o objetivo é transitar de  $\langle HLCR-0 \rangle$  para  $\langle 0-HLCR \rangle$ . Considerando que o bote suporta apenas o Homem e  mais um  item, descreva a sequência de símbolos (cadeia) que representa a solução segura (onde ninguém é comido).Exercício 12 (Diferença AFD vs AFND):A principal distinção entre AFDs e AFNDs reside na função de transição  $\delta$ . Utilizando a notação matemática, explique como o conceito de  conjunto das partes  (powerset ou  $2^S$ ) é aplicado na definição do AFND.Exercício 13 (Vending Machine):Imagine uma máquina que libera um produto ao acumular  30 centavos ou mais  ( $F = \{ \text{valores} \geq 30 \}$ ). Com  $\Sigma = \{5, 10, 25\}$ , trace o caminho dos estados ao inserir a sequência: 10, 5, 25. O estado final é atingido? Qual o valor total acumulado?Exercício 14 (Construção de AFD Simples):Construa um AFD sobre  $\Sigma = \{a, b\}$  que aceite apenas cadeias que iniciam com o símbolo 'a'. Preencha a tabela de transições abaixo (utilize  $S_{err}$  para o estado de erro):| $\delta$ | a | b || ------ | ------ | ------ || $S_0$  (Inic) |  |  || $S_1$  (Final) |  |  || $S_{err}$ |  |  |
Exercício 15 (Processamento de Cadeia):Considere o AFD recognitor da subpalavra "aa":
$S = \{S_0, S_1, S_2\}$ ,  $S_0$  inicial,  $F = \{S_2\}$ .
Transições:  $\delta(S_0, a)=S_1, \delta(S_0, b)=S_0; \delta(S_1, a)=S_2, \delta(S_1, b)=S_0; \delta(S_2, a)=S_2, \delta(S_2, b)=S_2$ .Demonstre o processamento passo a passo da cadeia abaa, indicando o estado ativo após cada símbolo.
3. Nível Difícil: Transformações, Minimização e Provas
Exercício 16 (Prova de Propriedade):Utilizando a definição formal de prefixo ( $x$  é prefixo de  $y$  sse  $\exists w \in \Sigma^$  tal que  $y=xw$ ), prove logicamente: Se  $x$  é prefixo de  $y$  e  $y$  é prefixo de  $z$ , então  $x$  é prefixo de  $z$ .Exercício 17 (Conversão AFND para AFD):No processo de determinização, um novo estado no AFD é criado a partir de um conjunto de estados do AFND. Explique como tratar o caso onde  $\delta(S_0, a) = \{S_0, S_1\}$ . O que o novo estado  $\{S_0, S_1\}$  representa em termos de processamento paralelo?Exercício 18 (Minimização):Dois estados são equivalentes se, para qualquer cadeia  $w$ , ambos levam a um estado final ou ambos levam a um estado não-final. Analise um caso onde  $S_3$  e  $S_4$  são finais e possuem as mesmas transições:  $\delta(S_3, a)=S_3, \delta(S_4, a)=S_3$  e  $\delta(S_3, b)=S_2, \delta(S_4, b)=S_2$ . Como a minimização simplifica esse cenário?Exercício 19 (Padrões Complexos):Desenhe a lógica de um AFD para  $L = \{w \in \{a, b\}^ \mid w \text{ possui quantidade par de 'a's e ímpar de 'b's} \}$ .Dica do Professor: Utilize quatro estados para representar o produto cartesiano das paridades:  $S_{pp}, S_{pi}, S_{ip}, S_{ii}$  .Exercício 20 (Análise de Sufixo):Construa a estratégia para um AFD que aceite cadeias onde o  quinto símbolo da direita para a esquerda  seja 'a'. Explique o conceito de "janela de memória" necessário para este autômato e por que o número de estados cresce exponencialmente ( $2^n$ ).
Gabarito Sugerido (Respostas Objetivas)
Ex 02:  a) 6; b) 4 (símbolos:  $a, b, de, c$ ); c) Não pertence (o símbolo  $d$  isolado não existe em  $\Sigma$ ).
Ex 04:   $\Sigma^*$  é infinito sempre que  $\Sigma \neq \emptyset$ .
Ex 06:  a)  $\{\lambda, a, ab, ba\}$ ; b)  $\{a\}$ .
Ex 08:   $\{00, 01, 000, 110, 111, 1100\}$ .
Ex 10:  Cadeias que terminam com o símbolo '0'.
Ex 11:  Cadeia:  $c, h, l, c, r, h, c$  (Homem atravessa com Cabra, volta sozinho, leva Lobo, volta com Cabra, leva Repolho, volta sozinho, leva Cabra).
Ex 12:  AFD:  $\delta: S \times \Sigma \to S$ ; AFND:  $\delta: S \times \Sigma \to \mathcal{P}(S)$ .
Ex 13:  Caminho:  $\langle 0c \rangle \xrightarrow{10} \langle 10c \rangle \xrightarrow{5} \langle 15c \rangle \xrightarrow{25} \langle 40c \rangle$ . Sim, estado final atingido. Total: 40c.
Ex 14:   $\delta(S_0, a)=S_1, \delta(S_0, b)=S_{err}, \delta(S_1, a/b)=S_1, \delta(S_{err}, a/b)=S_{err}$ .
Ex 15:   $S_0 \xrightarrow{a} S_1 \xrightarrow{b} S_0 \xrightarrow{a} S_1 \xrightarrow{a} S_2$  (Aceito).
Ex 18:   $S_3$  e  $S_4$  são redundantes e podem ser fundidos em um único estado.
Ex 20:  O autômato deve manter estados que representem todas as combinações possíveis dos últimos 5 símbolos lidos (necessários  $2^5 = 32$  estados).

\end{document}
